\section{Limitations}
\label{sec:threats}

\subsection{Construct validity}
LineVul fine-tunes its encoder end to end, whereas the principal graph results
use frozen features, so any absolute difference between them mixes
representation with training regime. We attempted to separate the two by training the encoder jointly with the graph network. Against a matched control, identical in data, schedule and architecture but with the encoder frozen, joint training was worth $-0.003$ (paired over five seeds, $p=0.26$). That measurement bounds the confound rather than removing it: on this pipeline, unfreezing the encoder is too small an effect to account for the gap between the strongest graph configuration ($0.658$) and LineVul ($0.792$) on the same partition. It does not follow that the two are comparable in other respects, and we do not claim the gap is therefore representational. The
internal ablations in Section~\ref{sec:graph-worth} are unaffected either way,
as the encoder is held fixed on both sides there.

\subsection{External validity}
All results are derived from a single C/C++ corpus. While the mechanism by which a permissive split admits near-copies of training data is not language-specific~\citep{allamanis2019adverse}, the measured magnitudes are specific to this context.

The decomposition is measured on one corpus and depends on how that corpus is assembled. How much of the leakage is attributable to duplication rather than to repository overlap is a function of how many projects the corpus spans and how heavily they are duplicated, so the proportions reported here should not be carried over to corpora with a different structure. The total, the overall gap between a permissive and a strict protocol, is the more portable quantity, but we measure it on one corpus and do not claim it is invariant.

\subsection{Encoding and architecture}
The edge-direction result holds for three of the four backbones and is
reversed in the fourth. At five seeds the graph transformer gains $+0.016$
($p=0.010$) and GCN $+0.006$ ($p=0.062$), while GIN loses $0.008$ ($p=0.043$).
The split follows whether the operator conditions on edge type, which is the
mechanism we claim, but three architectures is a narrow base for a general
recommendation. We also do not claim the reverse
encoding is optimal; alternatives such as basis-decomposed relational
convolutions~\citep{schlichtkrull2018modeling}, which would let a model with
per-relation weight matrices absorb eight relations more cheaply than four, are
untested here. Nor do we vary depth against the known cost of doing
so~\citep{oono2020graph,alon2021bottleneck}: the reverse-edge result widens the
receptive field at fixed depth, and trading depth against that width is a
separate experiment.

The positive control in Section~\ref{sec:graph-worth} establishes that the
pipeline learns labels defined by the edge set, using coarse global properties.
It does not establish that the same pipeline could localise a subtle dependency
change, and the paired evaluation suggests it cannot.

\subsection{What we do not claim}
No claim is made that graph representations are universally ineffective, nor that the evaluated architectures are optimally tuned. The claim is that, under the evaluation protocol and encoders commonly used in this literature, the measured contribution of the code property graph is $+0.003$ with the directed encoding the field uses, rising to $+0.032$ once reversed edges are restored, and that in either case the protocol accounts for a far larger share of reported performance than the representation does.

